% Primary source audit, completed 15 September 2026.
% 2015: PMC4409235 full primary HTML, Experimental Procedures (Statistical
% Analyses), Results, Table 1 and Figures 1--2; retrieved with urllib after
% the web reader returned a CAPTCHA. SHA256 of retrieved HTML:
% 72dd59caaede1c477c643f818926390bf48e2966692b4dbd47b752b4f6fbe98d.
% 2018: literature/core/2018_AE4_model.pdf, pp. 258--263, 269--271;
% SHA256 e2782a89b13a6dd27ea1c510d11515bacc33931697baccb678802de47bb13a23.
% 2021: primary PubMed record 34585968, abstract and Figures 1--2 captions;
% the complete article was not required or claimed to be reread.
% 2025: PMC12463136 primary author manuscript, Methods, Results and Discussion;
% retrieved HTML SHA256:
% 890c6f4b9834e8f0c8683e181aac5c8e3d49e045d7ada3d728b1b0cc591c2583.
% No source text or protected manuscript text is copied into this section.
\section{Comparison with experimental evidence}
\label{sec:literature-comparison}

\paragraph{The gland phenotype and the comparison threshold.}
In the 2015 mouse submandibular gland study, stimulation used
$0.3\,\mu\mathrm{M}$ carbachol and $5\,\mu\mathrm{M}$ isoproterenol.
AE4 deletion reduced total secretion over ten minutes by
$35\pm4.7\%$ (mean $\pm$ standard error; six glands per genotype).
Initial flow over two to three minutes was comparable, with the deficit
developing later. AE2 deletion had no detectable flow effect against its
own control strain. AE4 deficient acinar cells had lower resting chloride
and slower chloride reuptake; initial chloride exit and resting pH were
not detectably changed. These are distinct measurements, rather than
interchangeable manifestations of a single steady flux~\cite{Pena2015}.
The prescribed $30.3\%=35\%-4.7\%$ inverse target is therefore a
comparison value one reported standard error below the mean. It is not
a confidence bound, an uncertainty interval for a model parameter, or
a universal biological minimum.

The Task 40 cumulative null deficit of $3.8575\%$ establishes a small
effect for the frozen model and protocol. The equal routing identity
does not establish that small magnitude for every admissible parameter
choice. Nor does agreement at $600\,\mathrm{s}$ establish the observed
time course. The present simulations initialise every genotype at the
inherited WT resting state, whereas the experiment used established
knockout animals. That difference should remain explicit when comparing
the initial response or stored intracellular chloride.

\paragraph{Compensation and pH.}
The 2015 isolated NKCC1 assay used bicarbonate free solution,
$30\,\mu\mathrm{M}$ ethoxyzolamide and
$50\,\mu\mathrm{M}$ T16Ainh-A01; bumetanide identified the
NKCC1 dependent component. Initial chloride uptake showed no detectable
genotype difference. Stimulation induced NHE dependent alkalinisation
also showed no genotype difference~\cite{Pena2015}.
Neither finding measures the complete stimulated network integral.
Consequently, the approximately $23.16\%$ rise in integrated NKCC1
loading in the frozen Task 40 null is an unresolved experimental
tension. It must not be described as observed compensation. The local
equilibrium compensation derivative, the relative integrated NKCC1
increase and the fraction of missing AE4 chloride replaced by NKCC1
are three separate quantities throughout this report. Unchanged assay
alkalinisation does not establish equal NHE1 net flux at different ionic
states, and unchanged resting pH does not identify individual alkalinity
sources. A discriminating experiment would measure chloride loading,
pH, sodium and volume through the same combined stimulus, with
pathway specific perturbations and parallel controls.

\paragraph{What the 2018 model actually established.}
The published model predicted an approximately $24\%$ flow reduction
after AE4 deletion. Its reduced flow settled within about two minutes;
the authors explicitly contrasted this with experimental development
over roughly ten minutes. It predicted little NKCC1 activation following
either exchanger deletion. It included carbon dioxide hydration, NHE1,
AE2 and AE4, without another proton or bicarbonate transporter, and
held lumen volume constant~\cite{VeraSiguenza2018}.
Those are historical model properties, not experimental determinations
of the full network. The present retained carbon, alkalinity and volume
coordinates make the storage terms explicit. Numerical slow modes and
their projection onto secretion must be examined before attributing a
temporal discrepancy to absent slow regulation. A slow eigenvalue alone
does not reproduce delayed genotype divergence.

\paragraph{Regulation and the inverse CaCC family.}
The 2021 study found that beta adrenergic stimulation increased native
AE4 activity, that H89 prevented the response, and that constitutively
active PKA increased AE4 activity in CHO cells. The S173A mutant lost
PKA activation, whereas S273A retained it. The authors described
phosphorylation of S173 as a probable mechanism; H89 is a nonselective
inhibitor~\cite{Pena2021}.
This supports regulated AE4 activity, but does not measure the model's
AE4 regulatory gain, relaxation time or an AE4 to CaCC recruitment law.
The Task 41 inverse crossing quantifies a requirement within one family
selected using the phenotype. It supplies no independent validation of
that family. Its direct experimental test is the fraction of stimulated
apical chloride conductance lost upon AE4 perturbation at matched
calcium, membrane potential, chloride, pH and volume, resolved from the
initial response into sustained secretion. This distinguishes altered
channel recruitment from changed driving force and from intracellular
chloride storage.

\paragraph{Molecular source classes.}
Catal\'an and colleagues studied human AE4 expressed in HEK293 cells,
combining mutagenesis, transport assays and molecular dynamics. The
T756A--T448I mutant lost transport with extracellular sodium while
retaining activity with potassium. Their discussion proposes both
bicarbonate and carbonate cycles, including exchanged sodium and
potassium, and explicitly calls for further experiments to resolve
the alternatives~\cite{Catalan2025}.
Thus the seven frozen classes are specified source hypotheses, not seven
experimentally established full cycles. Their conserved projections
and equilibrium tests separate chloride, carbon, alkalinity and charge
consequences under inherited scalar kinetics. Matching one projected
signature cannot establish identical molecular mechanisms or identical
whole cell responses. An unsuccessful root search cannot establish
nonexistence. A thermodynamically opposed or nonphysiological case must
retain that qualification. Joint measurements of cation and anion
flux ratios, reversal conditions and current would constrain the source
vector; rates measured across those gradients would constrain the
separate kinetic law.

\paragraph{The most urgent missing measurements.}
The evidence checked here does not establish the identity, abundance
or capacity of the assumed inward $1\mathrm{Na}:2\mathrm{HCO}_3$
pathway. The alkalinity proposition gives a conditional requirement for
additional support when the admissible NHE1 and AE2 contribution is
insufficient. It does not prove that a particular NBC isoform is
universally necessary. The assumed loading allocation used to construct
NBC capacity is not a measured transporter partition.
Measurements of sodium and bicarbonate influx with current and pH
during combined stimulation are therefore a priority, alongside native
NHE1 and NKCC1 capacity and state dependence. The sensitivity crosswalk
identifies which effective capacities, conductances and buffer settings
matter most at the analysed point. Until independent ranges constrain
them, local elasticities do not establish robustness over biological
uncertainty.

\paragraph{Scope of source verification.}
The comparison uses the 2015 primary full text and Figures 1 and 2,
the repository copy of the 2018 published paper, the 2021 primary
abstract and figure captions, and the 2025 primary author manuscript
including its discussion. These four references were checked to that
extent. No systematic literature survey, new experimental measurement
or complete source audit of every inherited kinetic coefficient is
claimed here; parameter provenance is recorded separately in Task 43.
